\chapter{Metodologie simulačního modelu}
\label{chap:metodologie}

Tato kapitola popisuje metodologii stochastického simulačního modelu,
vyvinutého v~rámci této diplomové práce pro~hodnocení rizik investic
do~pěstování energetické biomasy (Miscanthus $\times$ giganteus a~SRC
vrby). Výklad sleduje logiku modelu od~volby lokality přes konfiguraci
parametrů až po~spouštění Monte Carlo simulace a~interpretaci výstupů.

Z~hlediska vnitřní mechaniky model generuje pro~každou plodinu $N$
nezávislých scénářů ročních peněžních toků po~celou dobu životnosti
plantáže, přičemž stochasticky modeluje (i)~výnos sušiny závislý
na~klimatu a~kvalitě půdy, (ii)~tržní cenu biomasy mean-revertním
Ornstein--Uhlenbeckovým procesem, (iii)~meziroční výkyvy počasí
včetně extrémních stresových let a~(iv)~riziko selhání plantáže
v~prvním roce. Z~výsledných distribucí peněžních toků se~vypočítají
agregované metriky výnosnosti (NPV, EAA, doba návratnosti, IRR)
a~rizika (směrodatná odchylka, Value at Risk, Conditional Value at Risk).

Vstupní parametry těchto sub-modelů jsou v~aplikaci předvyplněny
hodnotami kalibrovanými na~podmínky České republiky roku 2026;
uživatel je může před~spuštěním simulace upravit dle~konkrétního
investičního záměru. Podrobné odvození a~zdůvodnění jednotlivých
voleb následuje v~příslušných podsekcích této kapitoly.


%% ================================================================
\section{Lokalizace a~klimatická data}
\label{sec:lokalizace}

Prvním krokem uživatele je výběr lokality na interaktivní mapě. Kliknutím na mapu získá aplikace GPS souřadnice (zeměpisnou šířku $\varphi$ a~délku $\lambda$), ze kterých automaticky stahuje historická klimatická data z~API Open-Meteo \citep{OpenMeteo2023} z let 2014--2023: průměrnou denní teplotu a~roční úhrn srážek.

Na základě průměrné roční teploty $\bar{T}$, zeměpisné šířky a~ročních srážek je pozemek automaticky zařazen do jedné z~pěti klimatických zón:
\begin{itemize}
  \item \textbf{Tropická a~subtropická} ($|\varphi| < 23{,}5{}^{\circ}$ nebo $\bar{T} > 22\,{}^{\circ}$C),
  \item \textbf{Jižní Evropa / Středomoří} ($\bar{T} > 13\,{}^{\circ}$C),
  \item \textbf{Střední Evropa / mírné pásmo} ($\bar{T} > 8\,{}^{\circ}$C),
  \item \textbf{Severní Evropa / chladné pásmo} ($\bar{T} \leq 8\,{}^{\circ}$C),
  \item \textbf{Suché / marginální půdy} (roční srážky $< 450$\,mm).
\end{itemize}
Každá zóna má přiřazené výnosové rozsahy $[Y_{\min}, Y_{\max}]$ 
pro~obě plodiny, diferencované navíc podle kvality půdy. V~případě 
České republiky je kvalita půdy předvyplněna automaticky z~registru 
BPEJ (bonitovaná půdně ekologická jednotka) prostřednictvím API 
Výzkumného ústavu meliorací a~ochrany půdy (VÚMOP). Toto zařazení 
má však pouze doporučující charakter, neboť BPEJ je primárně určen 
pro konvenční zemědělské plodiny a~pro energetické plodiny se 
skutečná vhodnost může lišit. Uživatel by měl přiřazenou kategorii 
ověřit a~případně upravit, pokud disponuje přesnějšími údaji 
o~lokalitě. Mimo ČR volí uživatel kvalitu půdy ručně ze~čtyř 
kategorií: optimální, průměrná, neúrodná a~nevhodná.

Kompletní matice referenčních výnosových rozsahů pro všech pět 
klimatických zón a~čtyři kategorie kvality půdy je uvedena 
v~Tabulce~\ref{tab:vynosove-rozsahy}. Hodnoty pro střední Evropu 
byly stanoveny na~základě konzultací s~pracovníky VÚKOZ Průhonice 
z~dlouhodobých experimentálních ploch v~České republice 
a~jsou v~souladu s~publikovanými daty pro tento region 
\citep{mccalmont_2017,knapek_2024,stolarski_2019}. 

Hodnoty pro ostatní klimatické zóny představují \emph{doplňkovou 
část modelu} a~vycházejí výhradně z~mezinárodní odborné literatury, 
\textbf{bez konzultace s~odborníky pro dané regiony}. Jedná se 
proto o~orientační odhady na~základě literárních zdrojů, které 
slouží především k~ilustraci geografického rozšíření modelu 
a~před~použitím v~konkrétním investičním rozhodování v~zahraničí 
vyžadují validaci a~upřesnění regionálními experty.

Pro \textit{Miscanthus} $\times$ \textit{giganteus} v~teplejších 
zónách reflektují hodnoty vyšší výnosový potenciál C4~fotosyntézy, 
doložený globální analýzou modelu MiscanFor 
\citep{abdalla_2024}; konkrétně pro podmínky jižní Evropy 
vycházíme z~přehledové studie evropských polních pokusů 
\citep{LEWANDOWSKI2000}. Naopak v~severní Evropě jsou nižší 
hodnoty důsledkem citlivosti rhizomů na~zimní mrazy, dokumentované 
ve~tříletém pokusu v~Dánsku \citep{Jørgensen2003}, zatímco 
v~suchých a~marginálních podmínkách klesá výnos kvůli minimálnímu 
požadavku $\sim 450$~mm ročních srážek \citep{Ouattara2021}.

Pro RRD vrbu (rod \textit{Salix}) jsou rozsahy v~teplejších 
a~mírných klimatech v~souladu s~globálním datasetem výnosů 
\citep{castellanoalbors_2025}. V~severní Evropě hodnoty 
reflektují skutečnost, že vrba je domácím druhem skandinávského 
prostředí \citep{Verwijst13}, zatímco pro jižní Evropu vycházíme 
z~italské analýzy SRC plantáží \citep{BACENETTI2016} a~pro 
marginální půdy ze~studie \citep{Richard2019}.

Pod mapou aplikace zobrazí čtyři souhrnné metriky: průměrnou roční 
teplotu, roční úhrn srážek, detekované klimatické pásmo a~bonitní 
třídu půdy (BPEJ), čímž uživatel okamžitě vidí, v~jakém prostředí 
bude simulace probíhat.

\begin{table}[H]
\centering
\caption{Referenční výnosové rozsahy $[Y_{\min}, Y_{\max}]$ 
v~t~DM/ha podle klimatické zóny a~kvality půdy. }
\label{tab:vynosove-rozsahy}
\begin{tabular}{llcc}
\toprule
\textbf{Klimatická zóna} & \textbf{Kvalita půdy} 
& \textbf{Miscanthus} & \textbf{SRC vrba} \\
\midrule
Tropická / subtropická  & optimální & 20--30 & 18--25 \\
                        & průměrná  & 12--20 & 10--18 \\
                        & neúrodná  & 6--12  & 4--10  \\
                        & nevhodná  & 0--6   & 0--4      \\
\midrule
Jižní Evropa            & optimální & 18--28 & 12--18 \\
                        & průměrná  & 10--18 & 7--12  \\
                        & neúrodná  & 4--10  & 3--7   \\
                        & nevhodná  & 0--4   & 0--3      \\
\midrule
Střední Evropa          & optimální & 14--16 & 12--15 \\
                        & průměrná  & 7--14  & 7--12  \\
                        & neúrodná  & 4--7   & 4--7   \\
                        & nevhodná  & 0--4   & 0--4   \\
\midrule
Severní Evropa          & optimální & 8--12  & 6--10  \\
                        & průměrná  & 4--8   & 3--6   \\
                        & neúrodná  & 2--4   & 1--3   \\
                        & nevhodná  & 0--2   & 0--1      \\
\midrule
Suché / marginální      & optimální & 6--10  & 3--7   \\
                        & průměrná  & 3--6   & 1--4   \\
                        & neúrodná  & 1--3   & 0--1   \\
                        & nevhodná  & 0--1   & 0--1      \\
\bottomrule
\end{tabular}
\end{table}


%% ================================================================
\section{Konfigurace simulace}
\label{sec:konfigurace}

Ve druhém kroku uživatel nastavuje vstupní parametry simulace:
\begin{enumerate}
    \item \textbf{Výběr plodin} -- Miscanthus $\times$ giganteus, RRD vrba,
          nebo obě současně pro vzájemné srovnání.
    \item \textbf{Výměra plantáže} v hektarech (výchozí 10 ha), která slouží
          jako multiplikátor peněžních toků.
    \item \textbf{Počet Monte Carlo scénářů} $N$ (výchozí 10\,000, maximum 50\,000).
    \item \textbf{Podíl dotačního krytí} počáteční investice $s \in [0; 1]$.
    \item \textbf{Reálná diskontní sazba} $r$ pro výpočet čisté současné hodnoty
          (NPV); při 0~\% se pracuje s nominálním součtem.
    \item \textbf{Korelace $\rho$ mezi výnosovým a cenovým šokem}
          (výchozí $\rho = 0$, rozsah $[-1; +1]$). Bližší diskuse v sekci~\ref{sec:cenovy_model}.
    \item \textbf{Dekarbonizační drift $g$} (výchozí $g = 1{,}0~\%$/rok,
          rozsah $[-5; +5]~\%$/rok), reprezentující dlouhodobý reálný růst
          ceny biomasy. Bližší diskuse v sekci~\ref{sec:cenovy_model}.
    \item \textbf{Eskalace provozních nákladů $e_o$} (výchozí $e_o = 2{,}0~\%$/rok,
          rozsah $[0; 10]~\%$/rok), reprezentující roční růst nákladů na
          údržbu a sklizeň. Bližší diskuse v sekci~\ref{sec:naklady}.
    \item \textbf{Vlastnictví pozemku} (vlastní / pronajatý). Při~volbě
          pronajatého pozemku se~nastavuje pachtovné v~rozmezí
          120--280~EUR/ha/rok dle~bonity půdy.
    \item \textbf{Katastrofický scénář} (checkbox, default vypnuto)~---
          při zapnutí se~v~každém scénáři aplikuje souvislý desetiletý
          interval s~výnosem redukovaným na~40--60\,\% peaku
          (viz~sekce~\ref{sec:weather_mult}).

\end{enumerate}

%% ================================================================
\section{Nákladové parametry}
\label{sec:naklady}

Aplikace zobrazuje editovatelnou tabulku ekonomických parametrů pro
každou zvolenou plodinu. Uživatel může výchozí hodnoty upravit podle
vlastních podmínek. Výchozí parametry modelu jsou shrnuty v~Tabulce
\ref{tab:naklady-aktualizovane} a~kalibrovány na~české tržní podmínky
roku 2025 při kurzu 25~CZK/EUR.

Pro \textbf{SRC vrbu} máme k~dispozici aktuální a~podrobné podklady
přímo z~dlouhodobých experimentálních ploch VÚKOZ Průhonice
\citep{VUKOZ2026}, doplněné konzultacemi s~prof.~Knápkem
\citep{Knapek2026Komunikace}; jednotlivé položky CAPEX i~OPEX tak
vycházejí z~aktuální české terénní praxe.
U~\textbf{Miscanthusu} je situace podstatně horší~--- novější
ekonomické studie pro~podmínky ČR prakticky neexistují a~plodina
se~v~ČR komerčně pěstuje pouze ve~velmi omezeném rozsahu.
Jediným spolehlivě doloženým údajem je \textit{cena oddenků (rhizomů)}
od~německého dodavatele v~regionu Cheb~-- Františkovy Lázně, získaná
prostřednictvím doc.~Trogla \citep{Trogl2026}. Ostatní nákladové
položky (příprava pozemku, mechanizovaná výsadba, údržba, sklizeň,
likvidace) jsou \emph{expertně odhadnuty} na~základě (i)~staršího
metodického podkladu \citep{Strasil2009}, (ii)~analogie
s~odpovídajícími operacemi u~SRC vrby a~(iii)~dílčích konzultací
s~VÚKOZ. Konkrétní zdroj a~způsob odvození jednotlivých hodnot
je rozveden v~následujících podsekcích; tato datová asymetrie je
zároveň jedním z~limitů modelu diskutovaných v~kapitole
\ref{chap:zaver}.

\subsection*{Náklady na založení plantáže}

Náklady na založení plantáže představují dominantní nákladovou položku
celého životního cyklu obou plodin a~jsou v~modelu generovány
z~normálního rozdělení kolem výchozí hodnoty se~směrodatnou odchylkou
300~EUR (\textit{Miscanthus}) resp.~150~EUR (SRC), čímž je zachycena
variabilita mezi různými dodavateli sadby, kvalitou přípravy půdy
a~regionálními cenami práce.

U~\textit{Miscanthus} tvoří dominantní složku počátečních nákladů
sadební materiál~--- oddenky (rhizomy) v~hustotě 8\,000~ks/ha
\citep{Weger2021}. Cenová úroveň v~roce 2026 vychází ze~dvou zdrojů:
lokálního českého dodavatele v~regionu Cheb--Františkovy Lázně,
nabízejícího oddenky za~0,50~EUR/kus \citep{Trogl2026}, a~průměrné
nabídky evropských producentů v~rozsahu 0,16--0,21~EUR/kus
\citep{MiscanthusEU2026}. V~modelu uvažujeme kompromisní cenu
0,30~EUR/kus, která zahrnuje očekávanou dopravu z~EU dodavatele do~ČR
i~bezpečnostní rezervu na~manipulační ztráty. Celkové náklady
na~sadební materiál \textit{Miscanthus} tak činí přibližně 2\,400~EUR/ha.

Mechanizovaná výsadba probíhá u~\textit{Miscanthusu} adaptovanými
řízkovacími stroji, neboť oddenky lze do~půdy zakládat technologicky
obdobně jako řízky SRC vrby. Konkrétní agronomická data o~rychlosti
výsadby a~přepočtu na~jednotkovou sazbu nejsou pro~ČR k~dispozici;
\textbf{přebíráme proto expertní odhad} 1,5~Kč/oddenek převzatý
ze~sazby VÚKOZ Průhonice pro~výsadbu řízků SRC \citep{VUKOZ2026}.
Celkové náklady mechanizované výsadby \textit{Miscanthus} tak činí
$8\,000 \times 1{,}5~\text{K\v{c}} / 25 \approx 480$~EUR/ha.

U~SRC vrby je sadební materiál výrazně levnější: 7\,000~řízků/ha v~ceně
3~Kč/řízek (cca 840~EUR/ha) a~mechanizovaná výsadba ve~stejné struktuře
$7\,000 \times 1{,}5~\text{K\v{c}} / 25 \approx 420$~EUR/ha
\citep{VUKOZ2026}. 

Příprava pozemku zahrnuje u~obou plodin standardní operace~---
odplevelení, hlubokou podzimní orbu a~jarní vláčení~ a vychází na $\approx 104$~EUR/ha a~probíhá v~roce~0 (rok před~výsadbou). Vlastní výsadba
včetně sadebního materiálu probíhá v~roce~1, kdy je porost zároveň
vystaven nejvyššímu riziku selhání. Náklady na~manuální údržbu
v~prvních dvou letech (mechanické odplevelení) činí kumulativně
přibližně 145~EUR/ha pro~obě plodiny.

\subsection*{Roční provozní náklady}

Roční náklady na údržbu po fázi etablování se mezi plodinami liší.
U~SRC vrby jsou podle aktualizovaných expertních podkladů VÚKOZ
\citep{VUKOZ2026} u~dobře založené plantáže minimální a~zahrnují
především daň z~pozemku (cca 912~Kč/ha/rok, tj. 36~EUR/ha)
a~drobné provozní výdaje (mulčování meziřadí, kontrolní pochůzky),
celkem přibližně 50~EUR/ha/rok bez nájmu.
U~Miscanthus jsou roční náklady přibližně 100~EUR/ha/rok bez nájmu,
zahrnující především hnojení v~tříletém cyklu (60~kg~N/ha v~NPK formě)\citep{Weger2021}.

V~případě pronajatého pozemku se k~ročním nákladům připočítává
pachtovné, které se v~České republice typicky pohybuje v~rozsahu
3\,000--7\,000~Kč/ha/rok (120--280~EUR/ha/rok) v~závislosti na bonitě
půdy a~regionu \citep{VUKOZ2026}. Volba mezi vlastním
a~pronajatým pozemkem je v~modelu uživatelsky volitelná, neboť má
významný dopad na ekonomickou návratnost investice.

\subsection*{Náklady na sklizeň a~likvidaci plantáže}

Náklady na sklizeň jsou modelovány proporcionálně k~objemu sklizené
biomasy. U~\textit{Miscanthus} probíhá sklizeň každoročně v~předjarním
období při vlhkosti biomasy cca 20~\%. Pro sklizeň lze využít běžně
dostupné samojízdné \emph{řezačky určené primárně ke~sklizni siláže
kukuřice}, díky čemuž není pěstitel závislý na~úzce specializované
technice; jednotkové náklady tak \textbf{expertně odhadujeme}
na~přibližně 25~EUR/t, a~to analogií se~sazbami komerční
sklizně kukuřičné siláže v~ČR. U~SRC vrby probíhá sklizeň jednou
za~tři roky, přičemž omezená nabídka specializované sklizňové techniky
v~České republice~--- v~současnosti prakticky pouze jeden komerční
dodavatel \citep{VUKOZ2026}~--- vede k~vyšším jednotkovým nákladům
přibližně 40~EUR/t.

Likvidace plantáže \textit{Miscanthus} je technologicky jednoduchá~---
oddenky lze zničit standardní hlubokou orbou s~náklady přibližně
80~EUR/ha. U~SRC vrby je likvidace náročnější kvůli pařezům. Podle
aktualizované metodiky biologicko-mechanického rušení
\citep{VUKOZ2026} probíhá likvidace ve dvou až třech letech
po poslední sklizni, po postupném odumření pařezů (cca 2--3 roky) lze
pozemek zorat standardní agrotechnikou. Z~tohoto důvodu je u~SRC
vrby plánována poslední sklizeň v~26.~roce, biologicko-mechanické
rušení v~letech 27--28 a~konečná orba v~roce 29 tak, aby v~roce 30
mohl být pozemek vrácen do původního zemědělského využití. Celkové
náklady likvidace SRC vrby činí cca 6\,000~Kč/ha za orbu plus dvouletý
nájem během rušení (15\,450~Kč/ha celkem, tj. cca 620~EUR/ha).


\begin{table}[H]
\centering
\caption{Ekonomické parametry modelu (na 1~ha)}
\label{tab:naklady-aktualizovane}
\begin{tabular}{lrr}
\toprule
\textbf{Parametr} & \textbf{Miscanthus} & \textbf{SRC vrba} \\
\midrule
\multicolumn{3}{l}{\textit{Příprava pozemku a~výsadba (jednorázově)}} \\
\midrule
Příprava pozemku (EUR/ha)             & 104    & 104   \\
Sadební materiál (EUR/ha)             & 2\,400 & 840   \\
Mechanizovaná výsadba (EUR/ha)        & 480    & 420   \\
Údržba 1.--2.~rok (EUR/ha celkem)     & 145    & 145   \\
\textbf{Celkové CAPEX (EUR/ha)}       & \textbf{3\,129} & \textbf{1\,509} \\
\midrule
\multicolumn{3}{l}{\textit{Roční provozní náklady (od 3.~roku)}} \\
\midrule
Údržba bez~nájmu (EUR/ha/rok)         & 100      & 50       \\
Daň z~pozemku (EUR/ha/rok)            & 36       & 36       \\
Pachtovné~--- volitelné (EUR/ha/rok)  & 120--280 & 120--280 \\
\midrule
\multicolumn{3}{l}{\textit{Sklizeň a~likvidace}} \\
\midrule
Náklady sklizně (EUR/t)               & 25 & 40  \\
Likvidace plantáže (EUR/ha)           & 80 & 620 \\
\midrule
\multicolumn{3}{l}{\textit{Tržní a~strukturální parametry}} \\
\midrule
Výchozí prodejní cena (EUR/t~DM)      & 116  & 128  \\
Riziko selhání v~1.~roce              & 0{,}05 & 0{,}05 \\
Životnost plantáže (let)              & 25 & 29 \\
\bottomrule
\end{tabular}
\end{table}

%% ================================================================
\section{Výnosový model}
\label{sec:vynosovy-model}

Po konfiguraci parametrů a~stisknutí tlačítka \uv{Spustit simulaci} model nejprve pro každý scénář $i$ generuje výnosovou trajektorii. Maximální roční výnos $Y_{\max,i}$ je pro každý scénář generován z~normálního rozdělení $\mathcal{N}(\bar{Y},\, \sigma_Y^2)$, kde $\bar{Y}$ je střed výnosového rozsahu pro danou klimatickou zónu a~kvalitu půdy a~$\sigma_Y = (Y_{\max} - Y_{\min})/4$.

\subsection{\textit{Miscanthus $\times$ giganteus}}
\label{subsec:miscanthus-vynosy}

Pro \textit{Miscanthus $\times$ giganteus} je výnosový model kalibrován
na~empirickou výnosovou křivku publikovanou v~\citet{Weger2021}, která
vychází z~dlouhodobých polních pokusů VÚKOZ Průhonice (Michovky, Lukavec)
a~komerčních plantáží v~ČR. Křivka popisuje typický průběh ročního výnosu
sušiny v~průběhu životnosti plantáže~--- pomalou etablaci v~prvních
letech, prudký nárůst do~peaku ve~4.--5.~roce, plateau peakové
produktivity (rok 5--8) a~postupný úpadek od~roku~9 ke~konci ekonomické
životnosti.

Náš model přebírá tvar této křivky z~\citet{Weger2021} jako referenční
relativní profil, normalizovaný na~peak $=1{,}000$, a~aplikujeme jej
univerzálně přes všechny klimatické zóny a~kvality půdy~--- konkrétní
absolutní hodnota výnosu se~odvíjí od~maximálního ročního výnosu
$Y_{\max}$ pro~dané stanoviště. Konkrétní hodnoty relativního profilu
jsou shrnuty v~Tabulce~\ref{tab:miscanthus-yield-curve}.

\begin{table}[H]
\centering
\caption{Relativní výnosová křivka $\rho_M(t)$ Miscanthus $\times$ giganteus
v~průběhu životnosti plantáže . Hodnoty Y17--Y25 jsou extrapolací mírného pokračujícího úpadku za~horizontem
publikovaných dat.}
\label{tab:miscanthus-yield-curve}
\small
\begin{tabular}{cc|cc|cc}
\toprule
\textbf{Rok} & $\rho_M$ & \textbf{Rok} & $\rho_M$ & \textbf{Rok} & $\rho_M$ \\
\midrule
1 & 0{,}00 & 10 & 0{,}90 & 19 & 0{,}42 \\
2 & 0{,}03 & 11 & 0{,}77 & 20 & 0{,}40 \\
3 & 0{,}16 & 12 & 0{,}77 & 21 & 0{,}38 \\
4 & 0{,}94 & 13 & 0{,}71 & 22 & 0{,}36 \\
\textbf{5} & \textbf{1{,}00} & 14 & 0{,}65 & 23 & 0{,}35 \\
6 & 0{,}97 & 15 & 0{,}55 & 24 & 0{,}33 \\
\textbf{7} & \textbf{1{,}00} & 16 & 0{,}48 & 25 & 0{,}32 \\
\textbf{8} & \textbf{1{,}00} & 17 & 0{,}46 &  &  \\
9 & 0{,}94 & 18 & 0{,}44 &  &  \\
\bottomrule
\end{tabular}
\end{table}

\paragraph{Roční přírůstek sušiny.}
Roční výnos sušiny v~roce~$t$ pro Miscanthus plantáž s~maximálním ročním
výnosem~$Y_{\max}$ je dán přímým součinem relativního výnosu, maxima
a~stochastického počasového multiplikátoru:
\begin{equation}
    Y_{t}^{(M)} = \rho_M(t) \cdot Y_{\max} \cdot w_t,
    \label{eq:miscanthus-yield}
\end{equation}
kde $\rho_M(t)$ je hodnota z~Tabulky~\ref{tab:miscanthus-yield-curve},
$Y_{\max}$ je nominální maximální roční výnos sušiny dané kombinace
klimatické zóny a~kvality půdy (viz Tabulka~\ref{tab:vynosove-rozsahy}),
a~$w_t$ je stochastický počasový multiplikátor (viz § \ref{sec:weather_mult}).
Ve~špičkových letech ($\rho_M = 1$, roky 5, 7, 8) tedy roční výnos
odpovídá $Y_{\max} \cdot w_t$, v~ostatních letech je odpovídajícím
způsobem nižší. Sklizeň probíhá každoročně v~předjarním období při
vlhkosti biomasy cca 20~\%.

\paragraph{Charakteristické rysy křivky.}
Křivka má čtyři ekonomicky významné fáze:
\begin{enumerate}
    \item \textbf{Etablace} (rok 1--3): plantáž buduje rhizomový systém,
          výnos je zanedbatelný (0--16~\% peaku).
    \item \textbf{Skokový nárůst a~peak} (rok 4--8): rhizomový systém
          dospívá, výnos rychle dosahuje 94~\% peaku ve~4.~roce a~zůstává
          na~peakovém plateau (94--100~\%) do~8.~roku.
    \item \textbf{Postupný úpadek} (rok 9--16): pomalý pokles produktivity
          ze~$\sim$94~\% na~$\sim$48~\% peaku; plantáž zůstává
          ekonomicky životaschopná, ale s~klesající návratností.
    \item \textbf{Konec životnosti} (rok 17--25): mírný pokračující
          úpadek z~$\sim$46~\% na~$\sim$32~\% peaku; v~tomto stádiu
          již typicky probíhá rozhodování o~likvidaci plantáže
          a~návrat pozemku do~zemědělského využití.
\end{enumerate}


\subsection{SRC vrba -- výnosová křivka a~sklizňový cyklus}
\label{subsec:src-vynosy}

Pro rychle rostoucí dřeviny (SRC) je výnosový model odlišný od~Miscanthusu,
neboť biomasa je sklízena v~tříletých rotačních cyklech, nikoliv ročně.
Na~základě expertních dat VÚKOZ Průhonice \citep{VUKOZ2026} přiřazujeme
každé ze~čtyř kategorií kvality půdy (Nevhodná / Neúrodná / Průměrná /
Optimální) vlastní relativní výnosovou křivku, která popisuje, jakou
částí maximálního výnosu plantáže $Y_{\max}$ plantáž v~daném sklizňovém
cyklu skutečně dosahuje. Tvar křivky zachycuje typický průběh životnosti
plantáže~--- pomalý rozjezd v~prvních letech, peak ve~zralém stádiu a~postupný úpadek ke~konci životnosti, jehož rychlost závisí
na~kvalitě stanoviště. Konkrétní hodnoty pro~všechny čtyři kategorie
jsou shrnuty v~Tabulce~\ref{tab:src-yield-curves}.

\begin{table}[H]
\centering
\caption{Relativní výnos $\rho(t, q)$ SRC vrby pro 4 kvality půdy
(VÚKOZ Průhonice 2026). Hodnoty jsou relativní vůči maximálnímu výnosu
plantáže ($=1{,}000$); platí pro celý 3letý sklizňový cyklus s~uvedeným
rokem sklizně od~založení plantáže.}
\label{tab:src-yield-curves}
\small
\begin{tabular}{lcccccccc}
\toprule
\textbf{Kvalita půdy} & \multicolumn{8}{c}{\textbf{Rok sklizně od~založení plantáže}} \\
                                              & 5 & 8 & 11 & 14 & 17 & 20 & 23 & 26 \\
\midrule
Nevhodná   & 0{,}090 & 0{,}607 & 1{,}000 & 1{,}000 & 0{,}848 & 0{,}545 & 0{,}152 & 0{,}042 \\
Neúrodná   & 0{,}186 & 0{,}601 & 0{,}932 & 1{,}000 & 0{,}932 & 0{,}660 & 0{,}312 & 0{,}148 \\
Průměrná   & 0{,}355 & 0{,}768 & 0{,}955 & 1{,}000 & 0{,}916 & 0{,}749 & 0{,}497 & 0{,}329 \\
Optimální  & 0{,}309 & 0{,}877 & 1{,}000 & 0{,}974 & 0{,}859 & 0{,}686 & 0{,}500 & 0{,}364 \\
\bottomrule
\end{tabular}
\end{table}

\paragraph{Roční přírůstek sušiny.}
Výnos sušiny v~roce~$t$ pro plantáž s~maximálním ročním výnosem~$Y_{\max}$
a~kvalitou půdy~$q$ je dán přímým součinem relativního výnosu, maxima
a~stochastického počasového multiplikátoru:
\begin{equation}
    y_{\mathrm{SRC}}(t, q) = \rho(t, q) \cdot Y_{\max} \cdot w_t,
    \label{eq:src-yield-method}
\end{equation}
kde $\rho(t, q)$ je hodnota z~Tabulky~\ref{tab:src-yield-curves}
odpovídající aktuálnímu sklizňovému cyklu, $Y_{\max}$ je nominální
maximální roční výnos sušiny dané kombinace klimatické zóny a~kvality
půdy (viz Tabulka~\ref{tab:vynosove-rozsahy}), a~$w_t$ je stochastický
počasový multiplikátor (viz § \ref{sec:weather_mult}). Ve~špičkovém
roce plantáže ($\rho = 1$) tedy roční výnos odpovídá $Y_{\max} \cdot w_t$,
v~ostatních letech je odpovídajícím způsobem nižší.

\paragraph{Sklizňový cyklus.}
Sklizeň probíhá poprvé v~5.~roce projektu (rok výsadby $+\,4$~roky,
delší počáteční cyklus pro~etablování kořenového systému) a~následně
ve~tříletých intervalech (roky 8, 11, 14, \ldots, 26). Délka prvního
obmýtí pěti let byla potvrzena konzultací s~vědci VÚKOZ Průhonice
\citep{VUKOZ2026}. Sklizená biomasa v~roce sklizně je součtem ročních
přírůstků za~předchozí 3letý cyklus; v~letech bez sklizně plantáž
generuje pouze náklady na~údržbu.
\subsection{Stochastický meteorologický multiplikátor}
\label{sec:weather_mult}

Meziroční variabilitu výnosů způsobenou meteorologickými podmínkami zachycuje
multiplikátor $w_{t,i}$, který vstupuje do~výnosové
rovnice~\eqref{eq:miscanthus-yield} pro~Miscanthus a~analogicky do~vztahu
pro~roční přírůstek SRC vrby (rovnice~\eqref{eq:src-yield-method}).
Jde o~\textbf{bezrozměrnou relativní veličinu} generovanou z~normálního
rozdělení $\mathcal{N}(1{,}0;\, \sigma_w^2)$ s~ořezáním na~interval
$[0{,}5;\, 1{,}5]$ proti nerealistickým extrémům.

\paragraph{Význam parametru $\sigma_w$.}
Směrodatná odchylka $\sigma_w$ vyjadřuje typickou \textit{relativní}
fluktuaci výnosu mezi jednotlivými roky vůči deterministické střední
hodnotě dané růstovou křivkou. Hodnoty $w_{t,i} > 1$ reprezentují
nadprůměrně příznivý rok, hodnoty $w_{t,i} < 1$ rok podprůměrný.
Vzhledem k~normálnímu rozdělení s~jednotkovou střední hodnotou platí,
že přibližně 95~\% ročních realizací leží v~intervalu
$1 \pm 2\sigma_w$, což po~vynásobení očekávaným výnosem
$\bar{Y}$ odpovídá absolutnímu meziročnímu výkyvu o~přibližně
$\pm 2\sigma_w \cdot \bar{Y}$ tun sušiny na~hektar.

\paragraph{Kalibrace pro Miscanthus.}
Pro~Miscanthus je $\sigma_w = 0{,}15$, což odpovídá publikovanému
koeficientu variace meziročních výnosů Miscanthus v~podmínkách
České republiky ve~výši přibližně 15~\% \citep{knapek_2024}.

\paragraph{Kalibrace pro SRC vrbu.}
Pro~rychle rostoucí dřeviny je $\sigma_w = 0{,}10$, což odpovídá
empiricky pozorovanému koeficientu variace meziročních výnosů SRC
v~ČR ve~výši přibližně 10~\% \citep{knapek_2024}. Nižší variabilita
oproti Miscanthusu je dána (i)~robustnějším kořenovým systémem
dřevin tlumícím vliv krátkodobých meteorologických stresů
a~(ii)~víceletým akumulačním režimem mezi sklizněmi, který
průměruje meziroční výkyvy přírůstku.

\paragraph{Stresové roky.}
Kromě průběžné gaussovské variability model explicitně zahrnuje
\textit{extrémní meteorologické události} (stresový rok). S~pravděpodobností
$p_{\text{weather}} = 5\,\%$ (tj. průměrně 1~rok z~20) je výnos
v~daném roce dodatečně redukován na~50~\% očekávané hodnoty.
Tato pravděpodobnost a~míra poklesu byly stanoveny na~základě
expertního odhadu vědců VÚKOZ Průhonice a~odpovídají dokumentovaným
četnostem extrémních sucha v~ČR za~období 2000--2020
\citep{VUKOZ2026}.

\paragraph{Katastrofický scénář.}
Pro stresové testování dopadu dlouhodobě nepříznivého klimatu
(opakované sucho, vlna extrémních teplot) model nabízí volitelný
\emph{katastrofický scénář}, který uživatel aktivuje zaškrtnutím
v~konfiguračním kroku. Při zapnutí se~v~každém MC scénáři náhodně
vybere souvislý desetiletý interval v~produktivní fázi plantáže,
během něhož je výnosový multiplikátor v~každém roce nezávisle redukován
na~náhodný faktor v~rozsahu 40--60\,\%. Tento režim nahrazuje běžné
izolované stresové roky a~slouží k~vyhodnocení robustnosti projektu
vůči scénáři klimatické krize typu \uv{ztracená dekáda} s~vnitrodekádní
variabilitou intenzity sucha.


%% ================================================================
\section{Cenový model}
\label{sec:cenovy_model}

Vedle stochastické variability výnosů (sekce~\ref{sec:weather_mult})
představuje fluktuace tržní ceny biomasy druhý zásadní zdroj nejistoty
ovlivňující ekonomickou návratnost projektu. Vzhledem k~dlouhému
horizontu pěstování energetických plodin (25--29~let) má volba modelu
cenové dynamiky výrazný dopad na~výsledky simulace -- zejména
na~odhad pravděpodobnosti návratnosti, hodnoty Value at Risk
a~Conditional Value at Risk. Tato sekce nejprve formálně prezentuje diskretizovanou podobu modelu
s~parametry $\theta$, $\sigma_P$, $g$ a~$\mu_0$
(sekce~\ref{sec:ou_diskretizace}), a~v následujících podsekcích
se~věnuje stanovení \textbf{výchozí prodejní ceny} $P_0$
(sekce~\ref{sec:cena_kalibrace}) a~\textbf{kalibraci volatility}
$\sigma_P$ (sekce~\ref{sec:sigma_P}) na~základě
empirických dat, a~konečně \textbf{korelační struktuře} $\rho$
mezi výnosovými a~cenovými šoky (sekce~\ref{sec:korelace_drift}).

\subsection{Cenový model: diskretizovaný OU proces}
\label{sec:ou_diskretizace}

V~předkládaném modelu volíme pro~simulaci tržní ceny biomasy
\textbf{Ornstein--Uhlenbeckův mean-revertní proces}, který lépe
než geometrický Brownův pohyb odráží empirické chování komoditních
cen v~regulovaném prostředí (existence rovnovážné cenové hladiny,
substitučně-poptávková kontrola), viz. teoretická sekce~\ref{subsec:modelovani-cen}.

\paragraph{Diskretizace pro~roční krok.}
Pro~Monte Carlo simulaci s~ročním krokem ($\Delta t = 1$~rok) je
spojitý OU proces diskretizován pomocí \emph{Eulerova--Maruyamova
schématu}~--- standardní metody numerické integrace stochastických
diferenciálních rovnic, která v~každém časovém kroku aproximuje
deterministický drift a~stochastický šum nezávislými přírůstky
\citep{KloedenPlaten1992,Glasserman2003}. Diskretizovaná podoba má
tvar:
\begin{equation}
    P_{t+1,i} = P_{t,i} + \theta \cdot \bigl( \mu(t) - P_{t,i} \bigr)
                + \sigma_P \cdot \varepsilon_{t,i},
    \label{eq:ou_diskretni}
\end{equation}
kde $P_{t,i}$ je cena biomasy v~roce $t$ a~scénáři $i$ (EUR/t),
$\mu(t) = \mu_0 \cdot (1 + g)^t$ je časově rostoucí rovnovážná cena
a~$\varepsilon_{t,i}$ je standardně normovaný cenový šok.

\paragraph{Parametry modelu.}
Model vyžaduje pět parametrů, které jsou samostatně odvozeny
nebo kalibrovány v~následujících podsekcích:
\begin{itemize}
    \item \textbf{Výchozí cena} $P_0$~--- bottom-up odvození
          z~referenčního nákladu ERÚ a~transportu, viz~sekce~\ref{sec:cena_kalibrace}.
    \item \textbf{Volatilita} $\sigma_P$~--- empirická kalibrace
          pro~obě plodiny samostatně, viz~sekce~\ref{sec:sigma_P}.
    \item \textbf{Korelace} $\rho$ a~\textbf{drift} $g$~--- uživatelsky
          konfigurovatelné parametry, viz~sekce~\ref{sec:korelace_drift}.
\item \textbf{Rychlost návratu} $\theta = 0{,}25$~--- implikuje
      \emph{poločas návratu} ceny k~rovnovážné hladině
      $t_{1/2} = \ln(2)/\theta \approx 2{,}8$~roku. Hodnota leží
      uvnitř rozsahu empirických odhadů rychlosti mean-reverze
      pro~komoditní trhy reportovaných u~ropy a~mědi
      \citep{Schwartz1997} a~odráží regulatorní rámec lokálního
      českého trhu (cenové výměry ERÚ aktualizované v~ročním cyklu).
\item \textbf{Cenové bariéry} $[P_{\min};\, P_{\max}]$~--- slouží jako
      \textit{numerická pojistka} proti extrémním (a~ekonomicky
      nerealistickým) realizacím OU procesu v~dlouhém horizontu.
      Šíře intervalu je volena dostatečně velká vzhledem k~volatilitě
      $\sigma_P$ ($\sim\!\pm 3$--$5\,\sigma_P$ od~výchozí ceny),
      takže k~jejich aktivaci dochází jen výjimečně a~na~agregované
      metriky nemají podstatný vliv. Všechny hodnoty jsou shrnuty
      v~Tabulce~\ref{tab:cenovy_model_param}.
\end{itemize}

\paragraph{Strukturální dekompozice cenového šoku.}
Cenový šok $\varepsilon_{t,i}$ je strukturně rozdělen na~dvě nezávislé
složky: \textbf{70\,\% makroekonomická složka} (pocházející z~Gaussovské
kopule, viz~sekce~\ref{subsubsec:gauss-copula}, a~tedy potenciálně
korelovaná s~počasím přes parametr~$\rho$) a~\textbf{30\,\%
idiosynkratická složka} (čistě lokální, vždy nezávislá). Tento poměr
je pevný strukturální předpoklad modelu, nezávislý na~volbě~$\rho$;
zajišťuje, že cena obsahuje vlastní lokální šum (kurz, doprava,
krátkodobé výkyvy poptávky) i~v~případě silné korelace s~počasím.


% ============================================================
% Sekce 3.5 — Kalibrace ceny biomasy
% ============================================================

\subsection{Odvození ceny biomasy}
\label{sec:cena_princip}

Prodejní cena biomasy je v modelu vyjádřena v jednotkách \textbf{EUR za 
tunu sušiny} (EUR/t DM). Tato volba odpovídá výstupu výnosové funkce 
modelu, který přímo generuje hmotnost sušiny získanou ze sklizně 
(viz § \ref{sec:vynosovy-model}). Tržby v jednotlivých letech jsou pak prostým 
součinem výnosu sušiny a prodejní ceny:
\begin{equation}
R_t \;=\; Y_t \cdot P_t,
\label{eq:trzby}
\end{equation}
kde $Y_t$ [t DM/ha] je výnos sušiny v roce $t$ a $P_t$ [EUR/t DM] je 
prodejní cena.

Výchozí hodnota $P_0$ je odvozena vztahem
\begin{equation}
P_0 \;=\; P_\mathrm{ref}^{\mathrm{\check{C}R}} \cdot \mathrm{LHV}_\mathrm{DM} \;-\; T_\mathrm{DM},
\label{eq:cena_bottomup}
\end{equation}
kde $P_\mathrm{ref}^{\mathrm{\check{C}R}}$ je oficiální česká referenční 
cena biomasy [EUR/GJ], $\mathrm{LHV}_\mathrm{DM}$ je výhřevnost 
sušiny [GJ/t DM] a $T_\mathrm{DM}$ jsou náklady na svoz biomasy do 
teplárny přepočtené na 1 tunu sušiny [EUR/t DM]. Jednotlivé komponenty 
jsou kalibrovány v podsekcích \ref{sec:cena_pref}–\ref{sec:cena_kalibrace}.

% ------------------------------------------------------------
\subsection{Referenční cena biomasy $P_\mathrm{ref}^{\mathrm{\check{C}R}}$}
\label{sec:cena_pref}

Pro $P_\mathrm{ref}^{\mathrm{\check{C}R}}$ využíváme úředně stanovenou 
hodnotu nákladů na pořízení biomasy podle vyhlášky 
\citet{Vyhlaska315_2025} novelizující 
\citet{Vyhlaska79_2022}, příloha č. 2. Vyhláška stanovuje pro biomasu 
\textbf{kategorie 1 — cíleně pěstovaná pro energetické využití}, do níž 
spadají jak Miscanthus, tak rychle rostoucí dřeviny (RRD/SRC vrba), 
referenční cenu $190~\mathrm{K\check{c}/GJ}$ s účinností od 1.~1.~2026.

Tato hodnota je Energetickým regulačním úřadem (ERÚ) využívána pro 
stanovení měrných provozních nákladů tepláren a výpočet udržovací 
podpory tepla (viz \citet{ERU_CV8_2025}. Reprezentuje 
tedy cenu, za kterou je provozovateli teplárny ekonomicky dostupná 
vstupní biomasa kategorie~1.

Při stanoveném kurzu ($25{,}0$~CZK/EUR) odpovídá 
této hodnotě:
\begin{equation}
P_\mathrm{ref}^{\mathrm{\check{C}R}} \;=\; \frac{190}{25{,}0} \;\approx\; 7{,}60~\mathrm{EUR/GJ}.
\label{eq:pref_konverze}
\end{equation}


% ------------------------------------------------------------
\subsection{Výhřevnost sušiny $\mathrm{LHV}_\mathrm{DM}$}
\label{sec:cena_lhv}

Výhřevnosti sušiny obou plodin přebíráme z metodiky 
\citet{ThetaMetodika2024}, Tab.~4 (původní zdroj 
\citet{HavlickovaWegerKnapek2011}):
\begin{equation}
\mathrm{LHV}_\mathrm{DM}^{(M)} = 17{,}61~\mathrm{GJ/t~DM}, \qquad
\mathrm{LHV}_\mathrm{DM}^{(SRC)} = 18{,}39~\mathrm{GJ/t~DM}.
\label{eq:lhv_dm}
\end{equation}



% ------------------------------------------------------------
% ------------------------------------------------------------
\subsection{Náklady na svoz $T_\mathrm{DM}$}
\label{sec:cena_doprava}

Pro stanovení nákladů na silniční svoz biomasy do regionální výtopny 
kombinujeme dva nezávislé typy primárních dat: (i)~\emph{technologická data}
o objemu, hmotnosti a~časové náročnosti svozu z~konzultací 
\citep{VUKOZ2026}, a~(ii)~\emph{aktuální tržní sazby} dopravců pro 
nákladní soupravy nad~10~t užitečné hmotnosti dle 
\citep{AKOBlatny2025}.

\paragraph{Technologická data sklizně a~svozu.}
Podle expertních podkladů VÚKOZ Průhonice z~roku 2026 lze 
biomasu z~jednoho hektaru rotační sklizně SRC vrby (typický výnos 
$8\,$t\,DM/ha/rok, akumulovaný za~tříletý cyklus na~$24\,$t\,DM/ha) 
charakterizovat takto:
\begin{itemize}
    \item Hmotnost čerstvé štěpky: $\sim 50\,$t mokré (vlhkost~$\sim 52\,\%$),
    \item Objem čerstvé štěpky: $\sim 150\,$m$^3$ (sypná hmotnost~$\sim 333\,$kg/m$^3$),
    \item Počet potřebných souprav : 2,
    \item Celková doba svozu (obě soupravy, jízda + nakládka + vykládka): 5--6\,h.
\end{itemize}
Tato primární data slouží jako základ pro odvození měrné dopravní 
náročnosti, kterou v~modelu škálujeme na~zvolenou referenční vzdálenost.

\paragraph{Tržní sazby dopravců.}
Aktuální sazby pro nákladní soupravu o~hrubé hmotnosti $\sim 22$--$26\,$t
v~ČR jsou rozděleny na~kilometrickou složku (jízda) a~hodinovou složku
(nakládka a~vykládka). Z~veřejně publikovaného ceníku
\citep{AKOBlatny2025}, platného od~1.~března 2025, přebíráme:
\begin{align}
c_\mathrm{km}   &= 49\,\mathrm{K\check{c}/km} \;\;(\mathrm{bez\,DPH})
                = 1{,}96\,\mathrm{EUR/km}, \label{eq:sazba_km} \\
c_\mathrm{man}  &= 100\,\mathrm{K\check{c}/15\,min} \;\;(\mathrm{bez\,DPH})
                = 16\,\mathrm{EUR/h} \;\;(\mathrm{nakl.+vykl.}). \label{eq:sazba_man}
\end{align}

\paragraph{Náklad jedné jízdy.}
Při referenční svozové vzdálenosti $30\,$km jednosměrně (typická pro
lokální zásobování regionální výtopny / teplárny) a~jedné hodině
manipulace na~jízdu (kombinovaná nakládka + vykládka) má jedna jízda
náklad:
\begin{equation}
N_\mathrm{j\acute{\imath}zda} = 2 \cdot d \cdot c_\mathrm{km}
                                + t_\mathrm{man} \cdot c_\mathrm{man}
                              = 2 \cdot 30 \cdot 49 + 1 \cdot 400
                              = 3\,340\,\mathrm{K\check{c}}
                              \approx 134\,\mathrm{EUR/j\acute{\imath}zda}.
\label{eq:naklad_jizdy}
\end{equation}

\paragraph{Přepočet na 1\,t sušiny.}
Pro~kalibraci dopravních nákladů použijeme jako referenční výnos
průměr rozsahu \emph{Průměrná Střední Evropa}
z~Tab.~\ref{tab:vynosove-rozsahy}: $\bar{Y}^{(M)} = 10{,}5$~t~DM/ha/rok
pro~\textit{Miscanthus} a~$\bar{Y}^{(SRC)} = 9{,}5$~t~DM/ha/rok pro
SRC vrbu. Počet potřebných jízd na~hektar je u~SRC vrby určen
\emph{hmotnostním} limitem soupravy (25~t mokré štěpky), zatímco
u~Miscanthusu je limitujícím faktorem \emph{objemová} kapacita
(75~m$^3$) kvůli nízké sypné hmotnosti řezanky:
\begin{equation}
n_\mathrm{j\acute{\imath}zd}^{(SRC)} = \frac{M_\mathrm{ha}^{(SRC)}}{25\,\mathrm{t/souprava}}
\approx \frac{59{,}4}{25} \approx 2{,}4,
\qquad
n_\mathrm{j\acute{\imath}zd}^{(M)} = \frac{V_\mathrm{ha}^{(M)}}{75\,\mathrm{m}^3}
\approx \frac{105}{75} \approx 1{,}4,
\label{eq:n_jizd}
\end{equation}
kde $M_\mathrm{ha}^{(SRC)} = \bar{Y}^{(SRC)} \cdot 3 \cdot (1 - w_S)^{-1}
= 9{,}5 \cdot 3 / 0{,}48 \approx 59{,}4\,$t mokré štěpky akumulované
za~tříletý sklizňový cyklus při vlhkosti $w_S = 0{,}52$ (sypná
hmotnost čerstvé štěpky $\rho_\mathrm{syp}^{(S)} \approx 333\,$kg/m$^3$
\citep{Stolarski2019Chip}, který není vzhledem k~hmotnostnímu limitu
soupravy limitujícím faktorem), a~$V_\mathrm{ha}^{(M)} = \bar{Y}^{(M)}
\cdot (1 - w_M)^{-1} / \rho_\mathrm{syp}^{(M)} = 10{,}5 / 0{,}80 / 0{,}125
\approx 105\,$m$^3$ je objem řezanky z~jednoho hektaru Miscanthusu při~referenční
sypné hmotnosti $\rho_\mathrm{syp}^{(M)} = 125\,$kg/m$^3$ (střed publikovaného
rozsahu 100--140~kg/m$^3$ pro~Miscanthus chopped při~vlhkosti 15--20\,\%,
\citealp{Caslin2010,Smeets2009}) a~vlhkosti $w_M = 0{,}20$.

Měrné dopravní náklady přepočtené na~tunu sušiny:
\begin{align}
T_\mathrm{DM}^{(SRC)} &= \frac{n_\mathrm{j\acute{\imath}zd}^{(SRC)} \cdot N_\mathrm{j\acute{\imath}zda}}{\bar{Y}^{(SRC)} \cdot 3}
                       = \frac{2{,}4 \cdot 134}{9{,}5 \cdot 3}
                       \approx 11{,}3\,\mathrm{EUR/t\,DM},
\label{eq:T_DM_SRC} \\
T_\mathrm{DM}^{(M)}   &= \frac{n_\mathrm{j\acute{\imath}zd}^{(M)} \cdot N_\mathrm{j\acute{\imath}zda}}{\bar{Y}^{(M)}}
                       = \frac{1{,}4 \cdot 134}{10{,}5}
                       \approx 17{,}9\,\mathrm{EUR/t\,DM}.
\label{eq:T_DM_M}
\end{align}
Vyšší dopravní náklad u~Miscanthusu odráží jeho \emph{nižší sypnou
hmotnost}: souprava o~75~m$^3$ pojme pouze $\sim 9\,$t mokré řezanky
(oproti $\sim 25\,$t mokré štěpky SRC), což implikuje vyšší počet
jízd na~jednotku přepravené sušiny. SRC vrba má v~tomto ohledu
strukturální výhodu, neboť čerstvá dřevní štěpka má vyšší sypnou
hmotnost než drcená sláma. Per-jednotka sušiny je $T_\mathrm{DM}$
přitom \emph{invariantní} k~absolutní výši výnosu (vyšší $\bar{Y}$
implikuje úměrně více jízd a~úměrně více sušiny), záleží tedy
výhradně na~technologii dopravy a~fyzikálních vlastnostech materiálu.

\paragraph{Pozn.~k~diskrétnosti jízd.}
Vypočtené hodnoty $n_\mathrm{j\acute{\imath}zd}^{(M)} = 1{,}4$
a~$n_\mathrm{j\acute{\imath}zd}^{(SRC)} = 2{,}4$ představují
\emph{frakční amortizaci} přes typickou velikost komerční plantáže.
Pro modelovou plochu $\geq 10\,$ha je počet jízd $\lceil 10 \cdot
n_\mathrm{j\acute{\imath}zd} \rceil$ (tj.~14, resp.~24) prakticky rovný
součinu $10 \cdot n_\mathrm{j\acute{\imath}zd}$, takže frakční výpočet
odpovídá průměrné jednotkové hodnotě v~reálném provozu. Pro velmi
malou plantáž (< 5\,ha) by ceiling-efekt $T_\mathrm{DM}$ zvedl
o~cca $25$--$30\,\%$.

% ------------------------------------------------------------
\subsection{Kalibrované hodnoty $P_0$}
\label{sec:cena_kalibrace}

Dosazením do~vztahu~(\ref{eq:cena_bottomup}) získáváme:
\begin{align}
P_0^{(M)}   &= 7{,}60 \cdot 17{,}61 - 17{,}9  \;\approx\; 116\,\mathrm{EUR/t\,DM}, \\
P_0^{(SRC)} &= 7{,}60 \cdot 18{,}39 - 11{,}3  \;\approx\; 128\,\mathrm{EUR/t\,DM}.
\end{align}

Pro~kontextové srovnání uvádíme i~přepočet na~jednotku energie sušiny
po~odečtu dopravy:
\begin{equation}
\frac{P_0^{(M)}}{\mathrm{LHV}_\mathrm{DM}^{(M)}} = 6{,}59\,\mathrm{EUR/GJ}, \qquad
\frac{P_0^{(SRC)}}{\mathrm{LHV}_\mathrm{DM}^{(SRC)}} = 6{,}96\,\mathrm{EUR/GJ}.
\end{equation}
% ------------------------------------------------------------

% ============================================================
% Sekce 3.5.2 — Kalibrace volatility σ_P
% ============================================================
\subsection{Kalibrace volatility ceny biomasy}
\label{sec:sigma_P}

Volatilita $\sigma_P$ Ornstein--Uhlenbeckova procesu (vztah \ref{eq:ou-process}) je 
empiricky odvozena pro SRC vrbu z dlouhodobé veřejné cenové řady ČSÚ a 
následně přepočtena na Miscanthus pomocí mezinárodně dokumentovaného 
premia mezi bylinnou a dřevní biomasou.

% ------------------------------------------------------------
\paragraph{SRC vrba -- empirické odvození z ČSÚ.}
Pro rychle rostoucí dřeviny je k dispozici dlouhodobá veřejně dostupná 
cenová řada, kterou poskytuje Český statistický úřad 
\citep{CSU2025LesnictviI} v rámci publikace \emph{Indexy cen v lesnictví}. 
Jako proxy pro SRC vrbovou štěpku byla využita kategorie 
\emph{Dříví VI. třídy jakosti -- palivové dříví (vlastníci)}, představující 
listnaté dřevo nejnižší obchodní jakosti, energeticky využívané za obdobných 
podmínek jako SRC produkce. Časová řada zahrnuje 21 klouzavých čtvrtletních 
průměrů za období Q1/2019--Q1/2025.

Po konverzi z Kč/m$^3$ na EUR/t (předpoklad: hustota palivového dřeva 
$600\,$kg/m$^3$, výhřevnost $14{,}4\,$GJ/t, kurz $25\,$CZK/EUR) má časová 
řada průměrnou cenu $\bar{P}^{\mathrm{wet}} = 92{,}5\,$EUR/t, rovnovážnou 
směrodatnou odchylku $\sigma_\mathrm{eq}^{\mathrm{wet}} = 18{,}8\,$EUR/t a 
koeficient variace $\mathrm{CV} = 20{,}3\,\%$, s maximem v Q1/2023 
($121\,$EUR/t, vliv evropské energetické krize) a minimem v Q4/2020 
($73\,$EUR/t).

\paragraph{Přepočet na DM basis.}
Cenová řada ČSÚ je vyjádřena v EUR za tunu \emph{mokré} hmoty palivového 
dříví o vlhkosti $w \approx 20\,\%$ (odpovídá uvedené výhřevnosti 
$14{,}4\,$GJ/t). Jelikož v modelu pracujeme důsledně v jednotkách 
\emph{sušiny} (viz § \ref{sec:cena_princip}), provádíme přepočet všech 
statistik na DM basis vztahem
\begin{equation}
X^{\mathrm{DM}} \;=\; \frac{X^{\mathrm{wet}}}{1 - w},
\label{eq:dm_konverze}
\end{equation}
takže
\begin{align}
\bar{P}^{\mathrm{DM}}              &= 92{,}5 \,/\, 0{,}80 \approx 115{,}6 \;\mathrm{EUR/t\,DM}, \\
\sigma_\mathrm{eq}^{\mathrm{DM}}   &= 18{,}8 \,/\, 0{,}80 \approx 23{,}5 \;\mathrm{EUR/t\,DM}.
\end{align}


\paragraph{Odvození $\sigma_P$ z rovnovážné std.}
Pro odvození parametru $\sigma_P$ Ornstein--Uhlenbeckova procesu z 
rovnovážné směrodatné odchylky využíváme analytického vztahu platného 
v limitě dlouhodobé stacionarity:
\begin{equation}
\mathrm{Var}(P)_\mathrm{eq} \;=\; \frac{\sigma_P^2}{2\theta} 
\quad\Longleftrightarrow\quad 
\sigma_P \;=\; \sqrt{2\theta \cdot \mathrm{Var}(P)_\mathrm{eq}}.
\label{eq:OU_sigma_eq}
\end{equation}
Po dosazení $\theta = 0{,}25$ a 
$\mathrm{Var}(P)_\mathrm{eq}^{\mathrm{DM}} = 23{,}5^2 \approx 552\;\mathrm{EUR}^2/\mathrm{t\,DM}^2$ 
vychází $\sigma_P^{S,\mathrm{DM}} \approx 16{,}6\,$EUR/t\,DM. 
\paragraph{Miscanthus -- odhad na základě literatury.}
Pro Miscanthus nejsou v~ČR k~dispozici dlouhé časové řady tržních cen,
neboť jeho komercializace je v~rané fázi a~obchody se uskutečňují
převážně na~základě bilaterálních smluv mezi pěstitelem a~teplárnou, případně se neobchoduje vůbec.
Odhad volatility provádíme analogií: \citet{Kristoefel2014} empiricky
demonstrují na rakouských datech (měsíční ceny komodit, GARCH analýza),
že \emph{cenová volatilita zemědělské biomasy je v~průměru až několikanásobně
vyšší než dřevní biomasy}.

Konkrétní násobek volatility pro Miscanthus není v~literatuře přímo
kvantifikován; aplikujeme proto \emph{konzervativní přirážku}
$+20\,\%$ nad empirickou hodnotu pro SRC vrbu (dolní mez rozsahu z~literatury):
\begin{equation}
\sigma_P^{M,\mathrm{DM}} \;\approx\; 1{,}20 \cdot \sigma_P^{S,\mathrm{DM}}
\;=\; 1{,}20 \cdot 16,6 \;\approx\; 20\,\mathrm{EUR/t\,DM}.
\label{eq:sigma_M}
\end{equation}



\subsection{Korelace cenového a~výnosového šoku, dekarbonizační drift}
\label{sec:korelace_drift}

Dva parametry cenového modelu jsou v~aplikaci uživatelsky
konfigurovatelné: korelační koeficient $\rho$ mezi výnosovým
a~cenovým šokem a~dlouhodobý drift $g$ rovnovážné cenové úrovně.

\paragraph{Korelace $\rho$.}
V Gaussovské kopuli (viz § \ref{subsubsec:gauss-copula}) parametr $\rho \in [-1, +1]$ 
zavádí vazbu mezi výnosovým a cenovým šokem. V základní konfiguraci 
volíme $\rho = 0$, tj.~oba šoky jsou nezávislé. Důvodem je struktura 
českého trhu s energetickou biomasou: ten je objemově malý, regionálně 
fragmentovaný a slouží převážně k zásobování lokálních výtopen, takže 
lokální výkyvy nabídky nemají mechanismus, jak se promítnout do výkupní 
ceny --- ta je formována především centrálními faktory (ceny fosilních 
alternativ, regulace, dotace), nikoli tržním střetem nabídky a poptávky 
po vstupní biomase.

Záporná hodnota $\rho < 0$ by odrážela klasický komoditní mechanismus, 
kdy příznivé počasí zvyšuje produkci a rostoucí nabídka tlačí ceny dolů 
(a obráceně). Tento režim je relevantní spíše pro velké likvidní trhy 
a model jej dovoluje uživatelsky aktivovat nastavením $\rho < 0$ 
v parametrech, pokud by konkrétní tržní kontext nezávislost vyvracel. 
Citlivost výsledků modelu na volbu $\rho$ je analyzována v rámci OAT 
analýzy (§ \ref{sec:citlivost_oat}).

\paragraph{Dekarbonizační drift $g$.}
Drift $g$ rok určuje časový vývoj rovnovážného
cenového středu $\mu(t) = \mu_0 \cdot (1 + g)^t$.
Výchozí hodnota $g = 1{,}0\,\%$/rok odráží dlouhodobý reálný nárůst
ceny biomasy v~důsledku rostoucích cen emisních povolenek EU ETS
\citep{EuropeanCommission2024} a~obecných dekarbonizačních politik
podporujících substituci fosilních paliv obnovitelnými zdroji.
Tato hodnota je konzervativním odhadem, neboť projekce cen povolenek
pro~období 2026--2030 implikují růstové tempo v~rozmezí
1--3\,\%/rok \citep{ERU2025VymerBiomasa}.


\begin{table}[htbp]
\caption{Parametry cenového sub-modelu pro~obě plodiny ve~výchozí
konfiguraci. Hodnoty $P_0$ a~$\sigma_P$ odvozeny v~sekcích
\ref{sec:cena_kalibrace} a~\ref{sec:sigma_P};
$\theta$, $g$, $\rho$ a~bariéry sjednoceny pro~obě plodiny.}
\label{tab:cenovy_model_param}
\centering
\begin{tabular}{lcc}
\toprule
\textbf{Parametr} & \textbf{Miscanthus} & \textbf{SRC vrba} \\
\midrule
Výchozí cena $P_0$ (EUR/t DM)             & 116     & 128     \\
Volatilita $\sigma_P$ (EUR/t DM)          & 20      & 16{,}6  \\
Cenové bariéry $[P_{\min};\, P_{\max}]$ (EUR/t DM)
                                          & $[55;\, 200]$ & $[60;\, 220]$ \\
Rychlost návratu $\theta$ (--)            & 0{,}25  & 0{,}25  \\
Dekarbonizační drift $g$ (\%/rok)         & 1{,}0   & 1{,}0   \\
Korelace $\rho$ (--)                      & 0       & 0       \\
\bottomrule
\end{tabular}
\end{table}

%% ================================================================
\section{Selhání plantáže a~obnova}
\label{sec:selhani}

Model zahrnuje riziko selhání plantáže v~prvním roce s~pravděpodobností $p_f = 5\,\%$ (tj. přibližně 1~případ z~20). Tato hodnota byla stanovena na základě konzultací s~vědci VÚKOZ Průhonice. V~případě selhání jsou modelovány dva scénáře s~rovnou pravděpodobností:
\begin{itemize}
  \item \textbf{Obnova plantáže} (pravděpodobnost 50\,\%): pěstitel investuje dodatečných 60\,\% původního CAPEX do opakované výsadby a~plantáž pokračuje v~normální produkci od následujícího roku.
  \item \textbf{Opuštění projektu} (pravděpodobnost 50\,\%): investice je zcela ztracena, všechny budoucí peněžní toky jsou nulové.
\end{itemize}


%% ================================================================
%% ================================================================
%% ================================================================
\section{Peněžní toky}
\label{sec:cash_flow}

Pro~každý ze~$N$ scénářů Monte Carlo simulace generuje model kompletní
trajektorii ročních peněžních toků $\{CF_{t,i}\}_{t=0}^{T-1}$ po~celou
dobu životnosti projektu $T$. Struktura peněžního toku sleduje tři
chronologické fáze: nejprve \emph{investiční náklady} v~prvních třech
letech (roky 0--2), následně \emph{provozní fázi} se~stochastickými
tržbami ze~sklizně a~průběžnými náklady, a~v~závěru \emph{likvidaci
plantáže} v~posledních letech životnosti. Z~plné trajektorie
$\{CF_{t,i}\}$ jsou v~následující sekci~\ref{sec:vystupy} agregovány
finanční a~rizikové ukazatele prezentované uživateli.

\subsection{Investiční náklady a~obnova plantáže}
\label{subsec:cf_capex}

Investiční náklady (CAPEX) jsou rozloženy do~prvních tří let projektu
podle Tabulky~\ref{tab:naklady-aktualizovane}:
\begin{itemize}
    \item \textbf{ROK~0}: příprava pozemku (hluboká podzimní orba,
          jarní vláčení, odplevelení);
    \item \textbf{ROK~1}: výsadba (sadební materiál, mechanizovaná
          výsadba) a~první údržba; v~tomto roce je porost vystaven
          nejvyššímu riziku selhání;
    \item \textbf{ROK~2}: dokončení etablace (mechanické odplevelení,
          údržba 2.~roku).
\end{itemize}
Ve~vztazích~\eqref{eq:cf_misc} a~\eqref{eq:cf_src} jsou tyto náklady
promítnuty jako záporný příspěvek k~$CF_{t,i}$ v~příslušných letech.

V~roce výsadby (ROK~1) je peněžní tok dále ponížen o~případný náklad
na~\emph{obnovu plantáže} $\text{replant}_i$ s~podmíněnou pravděpodobností
$0{,}025$ (50\,\% selhání ve~5\,\%~scénářů, viz~sekce~\ref{sec:selhani}),
a~obě investiční složky (původní CAPEX i~obnova) jsou sníženy o~podíl
dotačního krytí $s$:
\begin{equation}
    CF_{1,i}^{\text{net}} = CF_{1,i} -
        \bigl(\text{CAPEX}_i + \text{replant}_i\bigr) \cdot (1-s).
    \label{eq:cf_year1_subsidy}
\end{equation}

\subsection{Provozní cash flow Miscanthus}
\label{subsec:cf_misc}

V~průběhu etablační fáze plantáž \textit{Miscanthus} produkuje pouze
zanedbatelnou biomasu: ROK~0--1 odpovídá přípravě a~výsadbě (žádný
výnos), v~ROK~2 dosahuje rhizomový systém přibližně 3\,\% peakové
produktivity (drobná raná biomasa). \textbf{První komerční sklizeň}
přichází ve~ROCE~3 (\textasciitilde 16\,\% peaku), prudký nárůst
v~ROCE~4 na~94\,\% peaku, plné plateau ROK~5--8 (viz~referenční
profil v~Tabulce~\ref{tab:miscanthus-yield-curve}). Sklizeň probíhá
každoročně v~předjarním období při~vlhkosti biomasy cca 20\,\%.

Pro~roky $t \geq 2$ má peněžní tok obecně tvar:
\begin{equation}
    CF_{t,i} = Y_{t,i} \cdot P_{t,i} - C_t^{\text{údržba}}
                - Y_{t,i} \cdot C_t^{\text{sklizeň}}
                - C^{\text{daň}} - C^{\text{pacht}},
    \label{eq:cf_misc}
\end{equation}
kde v~ROCE~2 je $C_t^{\text{údržba}}$ rovna polovině etablační údržby
($C^{\text{údržba}}_2 = C^{\text{ud,1--2}}/2$) a~výnos $Y_{2,i}
\approx 0{,}03 \cdot Y_{\max} \cdot w_2$ (raná biomasa); od~ROKU~3
je $C_t^{\text{údržba}}$ standardní roční údržba a~výnos sleduje
referenční profil $\rho_M(t)$.

\subsection{Provozní cash flow SRC vrby}
\label{subsec:cf_src}

Pro~SRC vrbu jsou tržby realizovány pouze ve~sklizňových letech
$\mathcal{T}_{\text{sklizeň}} = \{5, 8, 11, \ldots, 26\}$, zatímco
v~ostatních letech plantáž generuje pouze náklady na~údržbu mezi
sklizněmi. První sklizeň ve~5.~roce reflektuje delší počáteční cyklus
nutný pro~plné etablování kořenového systému dřeviny; následné rotační
sklizně probíhají v~tříletých intervalech až do~26.~roku:
\begin{equation}
    CF_{t,i} = \begin{cases}
        H_{t,i} \cdot P_{t,i} - C_t^{\text{údržba}}
        - H_{t,i} \cdot C_t^{\text{sklizeň}}
        - C^{\text{daň}} - C^{\text{pacht}},
        & t \in \mathcal{T}_{\text{sklizeň}}, \\[4pt]
        - C_t^{\text{údržba}} - C^{\text{daň}} - C^{\text{pacht}},
        & \text{jinak},
    \end{cases}
    \label{eq:cf_src}
\end{equation}
kde $H_{t,i}$ je akumulovaná sklizená biomasa dle~\eqref{eq:src-harvest}.

\subsection{Likvidace plantáže}
\label{subsec:cf_likvidace}

V~závěrečných letech životnosti projektu se~peněžní tok ponižuje
o~náklady na~likvidaci plantáže nutné pro~návrat pozemku
do~konvenčního zemědělského využití:
\begin{itemize}
    \item \textbf{\textit{Miscanthus}}: jednorázově 80~EUR/ha v~roce~25
          (snadné zničení oddenků hlubokou podzimní orbou);
    \item \textbf{SRC vrba}: celkem 620~EUR/ha rozložených do~let 27--29
          v~poměru $189 + 189 + 240$~EUR/ha dle~aktualizované metodiky
          biologicko-mechanického rušení (mulčování, postupný rozklad
          pařezů, finální orba) \citep{VUKOZ2026}.
\end{itemize}



%% ================================================================
\section{Výstupy modelu}
\label{sec:vystupy}

Výsledky simulace jsou v modelu strukturovány
do~sedmi tematických záložek umožňujících uživateli postupně procházet
ekonomickou výkonnost, výnosové a~cenové trajektorie, rizikový profil,
endogenně odvozenou diskontní sazbu, citlivostní analýzu, srovnání
kvalit půdy i~optimalizaci portfolia plodin. Každé záložce odpovídá
jedna z~následujících podsekcí; definice samotných ekonomických
a~rizikových metrik byly zavedeny v~teoretické
kapitole~\ref{chap:teorie-rizika-new} a~jsou zde aplikovány
na~simulační výstupy.

%% ----------------------------------------------------------------
\subsection{Záložka \uv{Přehled}}
\label{sec:vystupy_prehled}

První záložka shrnuje klíčové ekonomické ukazatele projektu.
V~horní části se~zobrazuje \emph{KPI karta} s~pěti hlavními metrikami
pro~každou plodinu: průměrný NPV (vztah~\eqref{eq:npv}), průměrná
ekvivalentní roční anuita EAA (vztah~\eqref{eq:eaa_theory}),
potenciál P95 (95.~percentil rozdělení NPV reprezentující optimistický
scénář), průměrná doba návratnosti $\overline{T^*}$ (vztah~\eqref{eq:payback})
a~empirická pravděpodobnost návratnosti $P(T^* < T_{\max})$.

Pod~KPI kartou následují dva grafické výstupy. \emph{Graf vývoje
peněžních toků} zobrazuje pro~každou plodinu průměrnou trajektorii
$\overline{CF_t}$ obklopenou pásmem 5.\,až 95.~percentilu, což názorně
ilustruje rozsah nejistoty v~jednotlivých letech. \emph{Histogram
celkového zisku} prezentuje empirické rozdělení $\{NPV_i\}_{i=1}^{N}$,
ze~kterého lze odečíst sešikmení rozdělení a~podíl ztrátových scénářů.

V~případě simultánního výpočtu obou plodin se~v~dolní části zobrazuje
\emph{rekapitulační tabulka} s~bočním srovnáním všech metrik. Z~této
záložky je~rovněž dostupný \emph{export do~formátu Excel} obsahující
sedm listů (souhrn, konfigurace, roční výnosy, ceny, náklady, cash
flow a~histogram NPV) pro~detailní analýzu mimo aplikaci.

%% ----------------------------------------------------------------
\subsection{Záložka Výnosy a~ceny}
\label{sec:vystupy_vynosy}

Tato záložka prezentuje stochastické trajektorie tří klíčových
veličin generovaných modelem. \emph{Trajektorie ročních výnosů
sušiny} zobrazuje pro~Miscanthus křivku $Y_{t,i}$ s~referenčním
piecewise-konstantním profilem dle~Tabulky~\ref{tab:miscanthus-yield-curve},
kolem kterého se~rozprostírá empirické pásmo P5--P95 odrážející
stochastickou variabilitu počasí. \emph{Sklizňový diagram SRC vrby}
ukazuje akumulovanou biomasu $H_t$ ve sklizňových letech, ilustrující
nárůst z nízkého výnosu v první rotaci k peakové produktivitě v rotacích
3–4 a následný úpadek (viz Tabulka \ref{tab:src-yield-curves}).

\emph{Vývoj tržní ceny} zobrazuje trajektorie OU procesu $P_{t,i}$
dle~vztahu~\eqref{eq:ou_diskretni} pro~obě plodiny, doplněné pásmem
spolehlivosti generovaným z~empirického rozdělení napříč scénáři.
Vizualizace umožňuje uživateli ověřit, zda použitý cenový model
reflektuje očekávanou dlouhodobou dynamiku včetně dekarbonizačního
driftu~$g$.

%% ----------------------------------------------------------------
\subsection{Záložka Riziko}
\label{sec:vystupy_riziko}

Záložka kvantifikuje riziko investice prostřednictvím tří doplňujících
se~ukazatelů odvozených z~empirického rozdělení $\{NPV_i\}_{i=1}^{N}$:
\begin{itemize}
    \item \textbf{Průměrná roční směrodatná odchylka} ročního cash flow,
          $\overline{\sigma(CF_i)}$, jako agregovaná míra meziroční
          nestability projektu;
    \item \textbf{Value at Risk na~hladině 5\,\%} odhadnutý jako
          5.~percentil empirického rozdělení NPV (viz~vztah~\eqref{eq:var-def});
    \item \textbf{Conditional Value at Risk na~hladině 5\,\%}
          odhadnutý jako průměr $NPV_i$ ve~scénářích padajících
          pod~hodnotu VaR (viz~vztah~\eqref{eq:cvar-def}).
\end{itemize}
Doplňkově je~zobrazen \emph{histogram rozdělení doby návratnosti}
$\{T^*_i\}$, který odhaluje pravděpodobnostní strukturu okamžiku
splacení investice.

\subsubsection{Dvourozměrná mřížková analýza}
\label{sec:citlivost_grid}

Pro~hlubší pochopení interakce dvou rizikových parametrů
záložka obsahuje \emph{teplotní mapu} (heatmap) systematicky variující
pravděpodobnost selhání plantáže $p_f \in [0;\, 0{,}20]$
a~pravděpodobnost extrémního klimatického stresu $p_w \in [0;\, 0{,}20]$
v~ekvidistantní mřížce $10 \times 10$. Pro~každou kombinaci $(p_f, p_w)$
je~provedena plná Monte Carlo simulace s~$N = 10\,000$ scénáři
a~vypočtena průměrná roční směrodatná odchylka cash flow. 

%% ----------------------------------------------------------------
\subsection{Záložka \uv{IRR a~diskont}}
\label{sec:vystupy_irr}

Tato záložka aplikuje \emph{endogenní odvození diskontní sazby}
z~výsledků modelu, jehož teoretický rámec byl zaveden v~podsekci
\ref{subsec:irr} kapitoly~\ref{chap:teorie-rizika-new}. Pro~každý
ze~$N$ scénářů Monte Carlo simulace se~řeší rovnice
$\mathrm{NPV}_i(r_i^*) = 0$ (vztah~\ref{eq:irr}), čímž vzniká
empirické rozdělení $\{r_i^*\}_{i=1}^{N}$ vnitřních výnosových
procent. Z~tohoto rozdělení se~následně odvozují tři varianty
doporučeného diskontu odpovídající různým preferencím investora
(rovnice~\ref{eq:irr_quantile} a~\ref{eq:irr_sharpe}).

\paragraph{Konfigurační formulář.}
V~horní části záložky se~nachází pre-filled formulář, v~němž uživatel
volí tři parametry odvození:
\begin{itemize}
    \item \textbf{Bezriziková sazba $r_f$} (default 4{,}0\,\%):
          odpovídá nominálnímu výnosu CZ 10Y státního dluhopisu
          k~květnu 2026, který představuje typickou bezrizikovou
          alternativu pro~investora v~České republice.
    \item \textbf{Cílový Sharpe ratio $S^*$} (default 0{,}4):
          typická hodnota pro~zemědělské investice; vstupuje
          do~Sharpe-based hurdle rate dle~vztahu~\eqref{eq:irr_sharpe}.
    \item \textbf{Akceptační kritérium $P(\mathrm{IRR} > r)$}
          (default 50\,\%, volitelně 75\,/\,90\,/\,95\,\%):
          pravděpodobnost, že IRR projektu překročí doporučený diskont.
          Volba odráží averzi konkrétního investora vůči riziku
\end{itemize}


\paragraph{Výstupy.}
Záložka prezentuje tři komplementární pohledy na~výsledné rozdělení:
\begin{itemize}
    \item \emph{Tabulka metrik IRR distribuce} (medián, průměr,
          směrodatná odchylka, uživatelsky zvolený akceptační percentil
          a~implikovaná riziková prémie $\pi = \mathbb{E}[\mathrm{IRR}] - r_f$);
    \item \emph{Tři varianty doporučeného diskontu} jako KPI karty:
          (A)~konzervativní kvantilový přístup
          (rovnice,
          (B)~Sharpe-based hurdle rate $r_f + S^* \cdot \sigma(\mathrm{IRR})$
          (rovnice,
          (C)~medián IRR jako break-even hodnota;
    \item \emph{Histogram IRR distribuce} s~vyznačenými úrovněmi
          všech tří doporučených diskontů, umožňující vizuální
          porovnání jejich umístění v~rámci empirického rozdělení.
\end{itemize}
Konkrétní hodnoty pro~obě plodiny napříč kvalitami půdy jsou
prezentovány v~kapitole výsledků~\ref{chap:vysledky}.
%% ----------------------------------------------------------------
\subsection{Záložka Citlivost}
\label{sec:vystupy_citlivost}

Záložka implementuje jednorozměrnou analýzu citlivosti metodou
\textit{one-at-a-time} (OAT), která představuje standardní lokální
přístup ke~kvantitativnímu hodnocení investičních projektů
\citep{Saltelli2008}.

\subsubsection{Parametry a~metodologie OAT}
\label{sec:citlivost_oat}

Uživatel volí z~dvanácti klíčových parametrů modelu
(Tabulka~\ref{tab:sa_params}); z~důvodu přehlednosti grafických
výstupů lze současně analyzovat nejvýše pět parametrů. Pro~každý
zvolený parametr je provedena série jedenácti simulací, ve~kterých
se~jeho hodnota systematicky mění podle pravidla~\eqref{eq:oat_param},
zatímco zbývající parametry zůstávají fixní na~výchozích hodnotách.

\begin{table}[htbp]
\caption{Parametry jednorozměrné analýzy citlivosti.}
\label{tab:sa_params}
\centering
\renewcommand{\arraystretch}{1.15}
\begin{tabular}{lll}
\toprule
\textbf{Kategorie} & \textbf{Parametr}  \\
\midrule
\multirow{5}{*}{Nákladové a~výnosové}
    & Prodejní cena (EUR/t)             \\
    & Výnosový potenciál (t/ha)          \\
    & Cena sadby (EUR/kus)             \\
    & Roční údržba (EUR/ha/rok)         \\
    & Náklady na~sklizeň (EUR/t)         \\
\midrule
\multirow{5}{*}{Rizikové a~strukturální}
    & Riziko selhání plantáže           \\
    & Životnost plantáže (let)          \\
    & Podíl dotačního krytí (\%)       \\
    & Diskontní sazba (\%)              \\
    & Prav.~klimatického stresu         \\
\midrule
\multirow{2}{*}{\shortstack[l]{Stochastické\\(absolutní osa)}}
    & Dekarbonizační drift $g$ (\%/rok)  \\
    & Eskalace nákladů $e_o$ (\%/rok)  \\
\bottomrule
\end{tabular}
\end{table}

\paragraph{Relativní vs.~absolutní osa.}
Pro~většinu parametrů je odchylka $\delta_k$ \textit{relativní}
v~rozsahu $\pm 10\,\%$ od~základní hodnoty s~krokem 2~p.b.
Dva parametry však používají \textit{absolutní} osu: dekarbonizační
drift $g$ ($[-5;\, +5]\,\%$/rok po~1~p.b.) a~eskalace nákladů $e_o$
($[0;\, 5]\,\%$/rok po~0,5~p.b.). Důvodem je, že ~relativní odchylka $\pm 10\,\%$ od~malé základní
hodnoty (typicky 1\,\%) by produkovala zanedbatelný rozsah variace.

\paragraph{Formální definice.}
Pro~ostatních deset parametrů s~vektorem výchozích hodnot
$\boldsymbol{\theta}^{(0)}$ a~relativní odchylkou
$\delta_k \in \{-0{,}10;\, -0{,}08;\, \ldots;\, +0{,}10\}$ je hodnota
$j$-tého parametru modifikována jako:
\begin{equation}
    \theta_j^{(k)} \;=\; \theta_j^{(0)} \cdot (1 + \delta_k),
    \qquad k = 1, 2, \ldots, 11,
    \label{eq:oat_param}
\end{equation}
zatímco ostatní parametry zůstávají na~výchozích hodnotách.
Pro~každou konfiguraci $\boldsymbol{\theta}^{(k)}$ je provedena
Monte Carlo simulace. 
\subsubsection{Sledované metriky a~vizualizace}

Z~výsledného rozdělení jsou pro~každou konfiguraci extrahovány
\emph{sedm ukazatelů}: (i)~průměrný NPV, (ii)~průměrná EAA
dle~\eqref{eq:eaa_theory}, (iii)~průměrná roční směrodatná odchylka
CF, (iv)~Value at Risk, (v)~Conditional Value at Risk
na~hladině $\alpha = 0{,}05$, (vi)~pravděpodobnost návratnosti
a~(vii)~průměrná doba návratnosti. Výsledky jsou vizualizovány jako
sedm spojnicových grafů, kde osa~$x$ představuje procentuální
(resp.~absolutní) odchylku parametru a~osa~$y$ hodnotu příslušné
metriky. Každý analyzovaný parametr je reprezentován samostatnou
barevnou křivkou, což umožňuje přímé porovnání citlivosti
a~identifikaci dominantních faktorů ekonomické výkonnosti.
Pod~spojnicovými grafy se~navíc zobrazuje \emph{tornádový graf}
sumarizující všechny OAT analýzy~--- parametry seřazené sestupně podle
absolutního dopadu na~NPV, čímž se~na~první pohled identifikují
nejvlivnější rizikové faktory.

%% ----------------------------------------------------------------
\subsection{Záložka Srovnání kvalit půdy}
\label{sec:vystupy_srovnani}

Záložka umožňuje paralelně spustit simulaci napříč všemi čtyřmi
kategoriemi kvality půdy zvolené klimatické zóny (Optimální, Průměrná,
Neúrodná, Nevhodná) a~prezentovat agregovaná srovnání. Hlavními výstupy
jsou:
\begin{itemize}
    \item \emph{Risk-return scatter graf}
          ($x$~=~směrodatná odchylka ročního cash flow,
          $y$~=~průměrná EAA), v~němž každá kombinace kvalita půdy
          $\times$ plodina představuje jeden bod, umožňující
          identifikovat efektivní hranici;
    \item \emph{Souhrnnou tabulku} s~klíčovými metrikami (EAA,
          $\sigma(CF)$, $P(\mathrm{NPV}>0)$, P5, P95) pro~všech
          osm kombinací (4 kvality $\times$ 2 plodiny);
    \item \emph{Procentuální vyjádření} výnosových ztrát při~přechodu
          z~optimální půdy na~horší kategorie, srozumitelně ilustrující
          dopad bonity na~ekonomickou návratnost.
\end{itemize}

%% ----------------------------------------------------------------
\subsection{Záložka Diverzifikace portfolia}
\label{sec:diverzifikace}

Závěrečná záložka hledá optimální poměr Miscanthus a~SRC vrby na~dané
ploše s~ohledem na~zvolený rizikový či~výnosový ukazatel. Návrh vychází
z~moderní teorie portfolia popsané v~sekci~\ref{subsec:system-nesystem-new}, dle~níž lze kombinací
nedokonale korelovaných aktiv snížit \textit{nesystematickou} složku
rizika oproti čisté monokultuře, aniž by~bylo nutné obětovat očekávaný
výnos. V~kontextu energetických plodin tento princip implikuje, že
kombinovaná plantáž obou plodin může vykazovat nižší volatilitu peněžních
toků (a~tedy lepší rizikové ukazatele VaR a~CVaR) než čistě monokulturní
řešení s~obdobnou očekávanou návratností.

Je třeba zdůraznit, že předkládaný optimalizátor představuje
\emph{zjednodušený} model: implicitně předpokládá nezávislé proporcionální
škálování CAPEX i~OPEX dle~zvoleného poměru ploch, a~neuvažuje praktické
dopady jako duplicitní investice do~specializované sklizňové techniky,
nutnost odděleného agronomického managementu, vyšší koordinační režii
či omezenou ekonomickou efektivitu malých dílčích bloků. V~reálném
provozu by skutečná diverzifikace tak vyžadovala nezanedbatelné dodatečné
náklady, které mohou potenciální rizikově-výnosový přínos kombinovaného
portfolia částečně eliminovat. 

\subsubsection{Optimalizační kritéria}
\label{sec:diverzifikace_metriky}

Uživatel volí jednu nebo více cílových metrik k~maximalizaci:
průměrný NPV, EAA, VaR 5\,\% nebo CVaR 5\,\%. Pro~každou zvolenou metriku $m$ je optimální
poměr definován jako
\begin{equation}
    w_M^*(m) = \arg\max_{w_M \in \{0, 2, \ldots, 100\}~\%}\, m\!\left(w_M\right).
    \label{eq:opt_mix}
\end{equation}
Vzhledem k~odlišnému dopadu metrik na~výsledek (NPV a~EAA typicky
upřednostňují plodinu s~vyšším výnosovým potenciálem, zatímco VaR
a~CVaR oceňují diverzifikační efekt) může být optimum pro~různé metriky
odlišné. Aplikace umožňuje souběžnou optimalizaci podle více metrik,
jejichž výsledky jsou prezentovány vedle sebe pro~snadné srovnání.

\subsubsection{Vizualizace výsledků}
\label{sec:diverzifikace_viz}

Pro každou vybranou metriku se~zobrazí:
\begin{itemize}
    \item KPI karta s~optimálním poměrem $w_M^* : (1-w_M^*)$, hodnotou
          metriky v~optimu a~rozdíly oproti čisté monokultuře M~/~SRC;
    \item Vizualizace pole (horizontální stacked-bar) ukazující
          rozdělení plochy mezi obě plodiny;
    \item Křivka \emph{efficient frontier} napříč všemi 51 poměry
          s~vyznačeným optimem (osa~$x$ = procento Miscanthus v~mixu,
          osa~$y$ = hodnota optimalizované metriky).
\end{itemize}

\section{Implementace}
\label{sec:implementace}

Model \textbf{BioFarm Simulator} je implementován v~jazyce \textbf{Python}
(verze 3.11+) s~uživatelským rozhraním postaveným na~frameworku
\textbf{Streamlit}; vektorizované numerické operace nad~maticemi scénářů
využívají knihovny NumPy a~SciPy, statistické rozdělení \texttt{scipy.stats},
interaktivní vizualizace knihovnu Plotly a~mapová komponenta integraci
\texttt{streamlit-folium}. Zdrojový kód je veřejně dostupný v~repozitáři
na~platformě \textbf{GitHub} pod~otevřenou licencí, což umožňuje
reprodukovatelnost simulací i~další rozšíření jinými výzkumníky.

\paragraph{Architektura aplikace.}
Uživatelské rozhraní je strukturováno do~\emph{čtyř očíslovaných
konfiguračních kroků} (lokalita a~klima, plodiny a~parametry, náklady,
výběr citlivostních parametrů) následovaných \emph{sedmi výsledkovými
záložkami} (Přehled, Výnosy a~ceny, Riziko, IRR a~diskont, Citlivost,
Srovnání kvalit půdy, Diverzifikace), z~nichž každá byla popsána
v~sekci~\ref{sec:vystupy}. Výsledky simulace jsou perzistovány
v~\texttt{st.session\_state}, takže přežijí změny ostatních widgetů
a~lze je opakovaně procházet napříč záložkami bez~nutnosti opětovného
spuštění simulace; náročnější výpočty (IRR distribuce, OAT citlivost,
diverzifikační optimalizátor) jsou triggovány samostatnými tlačítky
s~progress barem.

\paragraph{Výpočetní efektivita.}
Klíčové operace (Monte Carlo generace scénářů, OU diskretizace,
agregace metrik) jsou plně vektorizovány nad~maticemi tvaru
$(T \times N)$, čímž je dosaženo přijatelné doby výpočtu i~pro~$N = 10\,000$
scénářů na~běžném počítači.

%% ================================================================
